\chapter{电解质溶液}
我们已学习过一些有关电解质溶液和离子反应的初步知识，现在要运用物质结构和化学平衡等知识，来进一步学习电解质溶液的性质，更好地认识酸、碱、盐在水溶液里所起的反应，并了解原电池、电解、电镀等的化学原理。

\section{强电解质和弱电解质}
酸、碱和盐都是电解质，它们的水溶液都能导电。现在我们来进一步研究这样一个问题: 即在同样条件下，相同体积、 相同浓度而不同种类的酸、碱、盐的溶液，它们的导电能力是不是一样的。
\begin{Experiment}
按\cref{fig:5-1} 的装置把仪器连接好，然后把等体积的 \qty{0.5}{M} 的盐酸、氨水以及 \qty{0.5}{M} 醋酸、氢氧化钠、氯化钠的水溶液，分别倒入五个烧杯，连接电源。注意观察灯泡发光的明亮程度。
\tcblower
\begin{figurehere}
  \includegraphics{5-1.pdf}
  \caption{比较电解质溶液的导电能力}\label{fig:5-1}
\end{figurehere}
\end{Experiment}
实验结果表明: 连接插入醋酸溶液、氨水的电极上的灯泡比其它三个灯泡暗。
可见体积和浓度相同而种类不同的酸、碱和盐的水溶液在同样条件下的导电能力是不相同的。
盐酸和氢氧化钠、氯化钠溶液的导电能力比氨水和醋酸溶液强。

这是什么原因呢？
因为电解质溶液所以能够导电，是由于溶液里有能够自由移动的离子存在。
溶液导电性的强弱跟单位体积溶液里能自由移动的离子的多少有关。
也就是说，相同体积和浓度的导电性强的溶液里能自由移动的离子数目一定比导电性弱的溶液里的多。
这说明，电解质在溶液里电离的程度是不一样的。

我们知道，离子化合物是由阴离子和阳离子构成的。
如果我们把离子化合物氯化钠的晶体放入水中，在水分子的作用下，阴离子 \ce{Cl-} 和阳离子 \ce{Na+} 就会逐渐脱离晶体表面而进入溶液，成为能够自由移动的水合氯离子和水合钠离子。
在任何离子化合物的水溶液里，它们的阴、阳离子都同 \ce{Cl-} 和 \ce{Na+} 一样受水分子的作用，成为水合阴离子和水合阳离子。

为了简便起见，通常仍用普通离子的符号来表示水合离子。例如：
\begin{align*}
  \ce{NaCl}&=\ce{Na+ + Cl-}\\
  \ce{NaOH}&=\ce{Na+ + OH-}
\end{align*}
实验证明: 大多数盐类和强碱都是离子化合物，在它们的水溶液里，只有水合离子，没有分子。

具有极性键的共价化合物是以分子状态存在的。
例如，在液态氯化氢里只有氯化氢分子，没有离子存在。
氯化氢分子中，由于氯原子的电负性大于氢原子，氢原子跟氯原子结合的共价键是极性键。
氯化氢分子溶解于水时，在水分子的作用下，也能全部电离成为水合氢离子和水合氯离子，以致溶液里没有氯化氢分子存在。
其它的强酸如硫酸、硝酸等也跟氯化氢一样，它们的水溶液里只有水合氢离子和水合酸根离子存在。

由于氢离子是“裸露”的质子（氢原子失去电子后剩余由一个质子组成的核），半径很小，易被水分子吸引，生成水合氢离子，通常用 \ce{H3O+} 表示。
为了简便起见，我们也常把 \ce{H3O+} 写作 \ce{H+}。

但是某些具有极性键的共价化合物，它们溶解于水时，虽然同样受水分子的作用，却只有一部分分子电离成离子。
离子在互相碰撞时又互相吸引，而重新结合成分子。
因此，这种具有极性键的共价化合物在水里的电离过程是可逆的。
这个可逆的电离过程跟可逆的化学反应一样，它的相反的两种趋势，最终也将达到平衡。
因为在电离过程中，分子电离成离子的速度，必将随着溶液里离子的逐渐增多而减小，同时离子结合成分子的速度将不断增大。
在一定条件（如温度、浓度）下，两者的速度相等时，电离过程就达到了平衡状态，叫做\Concept{电离平衡}。
电离平衡跟化学平衡一样，也是动态平衡。
平衡时，单位时间里电离的分子数和离子重新结合生成的分子数相等，也就是说，溶液里离子的浓度和分子的浓度都保持不变。
所以这类具有极性键的共价化合物在水里的电离，常用可逆的电离方程式表示。
例如，醋酸和氨水的电离可以表示如下：
\begin{gather*}
  \ce{ CH3COOH <=> H+ + CH3COO-}\\
  \ce{ NH3.H2O <=> NH4+ + OH-}
\end{gather*}

由此可见，在这类电解质溶液里，既有离子存在，又有电解质分子存在。

根据上述三种情况可知: 离子化合物和某些具有极性键的共价化合物在水溶液里全部电离成为离子，没有分子存在。
所以不存在分子和离子之间的电离平衡，这样的电解质属于强电解质。
\Concept{强电解质}在水溶液里全部电离为离子。
如强酸、强碱和大部分盐类是强电解质。
某些具有极性键的共价化合物的水溶液里只有一部分电离成为离子，还有未电离的电解质分子存在，分子和离子之间存在着电离平衡。
这样的电解质属于弱电解质。
\Concept{弱电解质}在水溶液里只有部分电离为离子。
如弱酸、弱碱是弱电解质。

\begin{Practice}[习题]
\begin{question}
  \item 在氯化钠晶体里有没有离子存在？为什么氯化钠必须在水溶液里或熔化状态时才能导电？
  \item 在氯化氢的分子里有没有离子存在？为什么氯化氢的水溶液能够导电，而液态纯氯化氢不能导电呢？
  \item 在下列物质里哪些能够导电？为什么能够导电？写出电离方程式。哪些不能导电？为什么？
  \begin{tasks}(3)
    \task 氢氧化钾的水溶液，
    \task 氯化钾晶体，
    \task 醋酸的水溶液，
    \task 纯醋酸,
    \task 氯水，
    \task 液氯。
  \end{tasks}
  \item 在溶液导电性的实验装置里注入浓醋酸溶液时，灯光很暗，如果改用浓氨水，结果相同。可是把上述两种溶液混和起来实验时，灯光却十分明亮。为什么？
\end{question}
\end{Practice}


\section{电离度和电离常数}
\subsection{电离度}
不同的弱电解质在水溶液里的电离程度是不一样的。
有的电离程度大，有的电离程度小。
这种电离程度的大小，可用电离度来表示。
所谓电解质的\Concept{电离度}\emph{就是当弱电解质在溶液里达到电离平衡时，溶液中已经电离的电解质分子数占原来总分子数（包括已电离的和未电离的）的百分数}。
电解质的电离度常用符号 $\alpha$ 来表示：
\[ \alpha=\frac{\text{已电离的电解质分子数}}{\text{溶液中原有电解质的分子总数}}\times 100\% \]

例如，\qty{25}{\celsius} 时，在 $0.1M$ 的醋酸溶液里，每 \num{10000} 个醋酸分子里有 \num{132} 个分子电离成离子。
它的电离度是：
\[ \alpha=\frac{132}{\num{10000}}\times 100\% = 1.32\% \]

\cref{tab:5-1} 里是几种常见的弱电解质的电离度。
\begin{table}
  \caption{在 \qty{25}{\celsius} 时，$0.1M$ 溶液里某些弱电解质的电离度}\label{tab:5-1}
  \begin{tblr}{colspec={*2{cX[3,c]X[2,c]}},hline{2}=0.8pt,vline{4}=1.5pt}
    电解质 & 化学式     & 电离度（\%）& 电解质 & 化学式 & 电离度（\%）\\
    氢氟酸 & \ce{HF}    &  8.0 &   醋酸 & \ce{CH3COOH}  & 1.32 \\
    亚硝酸 & \ce{HNO2}  & 7.16 & 氢氰酸 & \ce{HCN}      & 0.01 \\
      甲酸 & \ce{HCOOH} & 4.24 &   氨水 & \ce{NH3.H2O}  & 1.33 \\
  \end{tblr}
\end{table}

从\cref{tab:5-1} 可见，在相同条件下，不同弱电解质的电离度不同，这是由弱电解质的相对强弱所决定的，一般说来，电解质越弱，电离度越小。所以电离度的大小，可以表示弱电解质的相对强弱。

电离度不仅跟电解质的本性有关，还跟溶液的浓度、温度等有关。
同一弱电解质，通常是溶液越稀，离子互相碰撞而结合成分子的机会越少，电离度就越大。
如在 \qty{25}{\celsius} 时，$0.2M$ \ce{CH3COOH} 的电离度为 0.948\%；$0.1M$ \ce{CH3COOH} 的电离度为 1.32\%；$0.001M$ \ce{CH3COOH} 的电离度为 13.2\%。
温度对电解质的电离度也有影响，当电解质分子电离成离子时，一般需要吸收热量，所以温度升高，平衡一般就向电离的方向移动，从而使电解质电离度增大。
因此，讲一种弱电解质的电离度时，应当指出该电解质溶液的浓度和温度。
如不注明温度，通常指 \qty{25}{\celsius}。

\subsection{电离常数}
弱电解质在一定条件下电离达到平衡时，溶液里各组分的浓度之间的关系，跟化学平衡一样，对一元弱酸或一元弱碱来说，溶液中电离所生成的各种离子浓度的乘积，跟溶液中未电离分子的浓度的比值是一个常数。
这个常数叫做\Concept{电离平衡常数}，简称为\Concept{电离常数}。
以醋酸为例，醋酸在水溶液里的电离方程式是：
\[ \ce{ CH3COOH <=> H+ + CH3COO- }\]
平衡时，溶液里各离子浓度的乘积，跟未电离分子的浓度的关系，可以下式表示：
\[ \frac{[\ce{CH3COO-}][\ce{H+}]}{[\ce{CH3COOH}]}=K_\text{电离}\]
式中 $[\ce{CH3COO-}]$ 和 $[\ce{H+}]$ 分别表示溶液中醋酸根离子的摩尔浓度和氢离子的摩尔浓度，$[\ce{CH3COOH}]$ 表示未电离的醋酸分子的摩尔浓度，$K_\text{电离}$ 表示醋酸的电离常数。

从醋酸的电离常数的表达式可以看出：$K_\text{电离}$ 值大，离子浓度必然也大，即表明该电解质较易电离。
所以，从 $K_\text{电离}$ 值的大小也可以看出弱电解质的相对强弱。
例如，醋酸和氢氰酸都是弱酸，已知 \qty{25}{\celsius} 时 $0.1M$ 醋酸溶液中醋酸的 $K_\text{电离}$ 值是 \num{1.75e-5}，而 $0.1M$ 氢氰酸溶液中氢氰酸的 $K_\text{电离}$ 值是 \num{4.93e-10}。
所以，氢氰酸是比醋酸更弱的酸。

对于同一弱电解质的稀溶液来说，电离常数跟化学平衡常数一样，不随浓度的变化而变化，而只随温度的变化而变化。
如在 \qty{25}{\celsius} 时醋酸的 $K_\text{电离}$ 是 \num{1.75e-5}，在 \qty{0}{\celsius} 时是 \num{1.65e-5}。
由于电离常数随温度的变化不大，在室温时，可以不考虑温度对电离常数的影响。

下面简单讨论多元弱酸、弱碱的电离。

多元弱酸的电离是分步进行的，例如：
\begin{gather*}
  \ce{ H3PO4 <=> H+ + H2PO4- }\\
  \ce{ H2PO4- <=> H+ + {HPO4}^2- }\\
  \ce{ {HPO4}^2- <=> H+ + {PO4}^3- }
\end{gather*}

它的每一步电离都各有电离常数，这些电离常数也各不相同。 
通常用 $K_1$、$K_2$、$K_3$ 等来加以区别。
不同的电解质具有不同的电离常数。
\cref{tab:5-2} 里是几种常见弱电解质的电离常数。
\begin{table}
  \caption{常见的几种弱电解质的电离常数（\qty{25}{\celsius}）}\label{tab:5-2}
  \begin{tblr}{colspec={*2{X[c]X[2,c]}},hline{2}=0.8pt,vline{3}=1.5pt}
    电解质 & 电离常数 & 电解质 & 电离常数 \\
    醋酸 \ce{CH3COOH} & \num{1.75e-5} & {磷酸 \\ \ce{H3PO4}} & {$K_1=\num{7.52e-3}$ \\ $K_2=\num{6.23e-8}$ \\ $K_3=\num{2.2e-13}$} \\
    {碳酸 \\ \ce{H2CO3}}   & {$K_1=\num{4.3e-7}$ \\ $K_2=\num{5.6e-11}$} & 亚硫酸 \ce{H2SO3} & {  $K_1=\num{1.54e-2}$（\qty{18}{\celsius}）\\ $K_2=\num{1.02e-7}$（\qty{18}{\celsius}）} \\
    {氢氰酸 \\ \ce{HCN}}   & $K=\num{4.93e-10}$ & {氢硫酸 \\ \ce{H2S}} & {  $K_1=\num{9.1e-3}$（\qty{18}{\celsius}）\\ $K_2=\num{1.1e-12}$（\qty{18}{\celsius}）} \\
    {氢氟酸 \\ \ce{HF}}   & $K=\num{7.2e-4}$ & {氨水 \\ \ce{NH3.H2O}} & \num{1.77e-5} \\
  \end{tblr}
\end{table}

比较多元弱酸各步电离的电离常数，可以看出，$K_1>K_2>K_3$。
以磷酸为例，$K_1$ 比 $K_2$ 约大 \num{e5} 倍，$K_2$ 比 $K_3$ 约大 \num{e5} 倍。
可见多元弱酸溶液的酸性主要由第一步电离所决定。

一元弱碱和多元弱碱的电离跟一元弱酸和多元弱酸的电离情况是相似的。

前面已经提到电离常数和电离度都可以表示弱电解质的相对强弱。
那么它们之间究竟有什么关系呢？让我们以 \ce{CH3COOH} 为例来说明。

设醋酸溶液里醋酸分子的摩尔浓度为 $c$，醋酸的电离度为 $\alpha$，那么溶液中每有 $c\alpha$\,\unit{mol/L} 的醋酸电离，就有 $c\alpha$\,\unit{mol/L} \ce{H+} 离子和 $c\alpha$\,\unit{mol/L} \ce{CH3COO-} 离子生成,即：
\[ [\ce{H+}]=[\ce{CH3COO-}]=c\alpha \]
醋酸的电离方程式是：

\begin{tblr}{colspec={*5{c}},hlines=0pt,vlines=0pt}
  & \ce{CH3COOH} & \ce{ <=> } & \ce{H+} & \ce{ + CH3COO-} \\
  开始浓度 & $c$ & & 0 & 0 \\
  开始浓度 & $c-\alpha$ & & $c\alpha$ & $c\alpha$ \\
\end{tblr}


\[ K_\text{电离}=\frac{[\ce{H+}][\ce{CH3COO-}]}{[\ce{CH3COOH}]}=\frac{c\alpha\cdot c\alpha}{c-c\alpha}=\frac{c\alpha^2}{1-\alpha}\]

对弱电解质来说，当 $K_\text{电离}$ 很小（$K_\text{电离}<\num{e-4}$）时，$\alpha$ 值也很小，所以可近似地认为 $1-\alpha\approx 1$。于是
\begin{gather*}
  K_\text{电离}=\frac{c\alpha^2}{1-\alpha}=c\alpha^2\\
  \alpha=\sqrt{\frac{K_\text{电离}}{c}}
\end{gather*}

这个公式表明: 在一定温度下,当溶液的浓度 $c$ 改变时，电离度 $\alpha$ 也随着改变，浓度越大，电离度越小。
我们已经知道，电离常数 $K_\text{电离}$ 是不随浓度改变而改变的。

由此可见电离常数比电离度能更好地表示出弱电解质的相对强弱。

上面已经知道电离度跟电离常数的关系，现在我们来运用它进行一些简单的计算。

\begin{example}
  在 \qty{25}{\celsius} 时，$0.10M$ 醋酸溶液中氢离子的浓度是多少？
\end{example}
\begin{solution}
  已知 $c=0.10M$，$K_\text{电离}=\num{1.75e-5}$
  \[ \alpha=\sqrt{\frac{K_\text{电离}}{c}}=\sqrt{\frac{\num{1.75e-5}}{0.10}}=\num{1.32e-2}\]
  $[\ce{H+}]=c\alpha=0.1\times\num{1.32e-2}=\num{1.32e-3}{M}$

  答: $0.1M$ 醋酸溶液中氢离子浓度为 $\num{1.32e-3}{M}$。
\end{solution}

\begin{example}
  在某温度时，已知 $0.20M$ 氨水的电离度是 0.934\%，求氨水的电离常数。
\end{example}
\begin{solution}
  已知 $c=0.20M$，$\alpha=0.934\%=0.00934$，代入
  \[ \alpha=\sqrt{\frac{K_\text{电离}}{c}}\]
  $K_\text{电离}=0.20\times(0.00934)^2=\num{1.74e-5}$

  答：在某温度时，氨水的电离常数是 \num{1.74e-5}。
\end{solution}

\begin{Practice}[习题]
\begin{question}
  \item 比较醋酸、氢氰酸、氢氟酸这三种酸的相对强弱。
  \item 为什么弱电解质的电离度跟溶液的浓度有关？
  \item 在 \qty{1}{L} \qty{2}{mol} 的电解质溶液里，有 \qty{0.2}{mol} 的电解质电离成离子。问这种电解质的电离度是多少？
  \item 在氟化氢溶液中，已电离的氟化氢为 \qty{0.2}{mol}，未电离的氟化氢为 \qty{1.8}{mol}。求该溶液中氟化氢的电离度。
  \item 某种电解质电离度大，是否它的溶液的导电能力就一定强？
  \item 下列两种溶液中，哪一种溶液的氢离子浓度大？为什么？
  \begin{tasks}
    \task \qty{1}{L} $0.1M$ 的醋酸溶液，
    \task \qty{1}{L} $0.01M$ 的醋酸溶液。
  \end{tasks}
  \item 在 \qty{25}{\celsius} 时，氢氰酸的 $K_\text{电离}$ 为 \num{4.93e-10}。计算 $0.01M$ 氢氰酸的电离度是多少？
  \item $0.1M$ 醋酸溶液的电离度是 1.32\%。求醋酸的 $K_\text{电离}$。
  \item 碳酸是弱酸，写出它的电离方程式，并指出碳酸在溶液里有哪几种离子？哪种最多？哪种最少？
\end{question}
\end{Practice}

\section{水的电离和溶液的 pH 值}
研究电解质溶液时往往涉及溶液的酸碱性。
电解质溶液的酸碱性跟水的电离有着密切的关系。
为了从本质上认识溶液的酸碱性，就要研究水的电离情况。

\subsection{水的电离}
根据精确的实验证明，水是一种极弱的电解质，它能微弱地电离，生成 \ce{H3O+} 和 \ce{OH-}。
\[ \ce{H2O + H2O <=> H3O+ +OH-} \]
可简写为：
\[ \ce{ H2O <=> H+ + OH- } \] 
\begin{figure}
  \includegraphics{5-2.pdf}
  \caption{水分子电离过程示意图}\label{fig:5-2}
\end{figure}

从纯水的导电实验测得，在 \qty{25}{\celsius} 时，纯水中 \ce{H+} 和 \ce{OH-} 的浓度各等于 \qty{e-7}{mol/L}。
在一定温度下，它跟其它弱电解质一样，有一个电离常数。
\[ K_\text{电离} = \frac{[\ce{H+}][\ce{OH-}]}{[\ce{H2O}]}\]

由于水的电离度很小，含有 \qty{55.5}{mol} 水分子的一升水中，仅有 \qty{e-7}{mol} 的水分子电离，它的已电离部分可以忽略不计，所以电离前后，水分子的摩尔数几乎不变，可以看做是个常数，因而上式可以写为：
\[ [\ce{H+}][\ce{OH-}]=K_\text{电离}[\ce{H2O}]\]
既然 $K_\text{电离}$ 是常数，$[\ce{H2O}]$ 也可以看做是个常数，常数乘常数必然成为一个新的常数，通常我们把它写作 $K_w$。即
\[ [\ce{H+}][\ce{OH-}]=K_w\]

$K_w$ 是水中 $[\ce{H+}][\ce{OH-}]$ 的乘积。
我们把 $K_w$ 叫做\Concept{水的离子积常数}，简称为\Concept{水的离子积}。 
$K_w$ 的值是多少呢？
已知 \qty{25}{\celsius} 时水中 \ce{H+} 和 \ce{OH-} 的浓度都是 \qty{1e-7}{mol/L}。所以
\[ K_w=[\ce{H+}][\ce{OH-}]=\num{1e-7}\times\num{1e-7}-\num{1e-14}\]

因为水的电离过程是一个吸热过程，所以当温度升高时，水的电离度增加，离子积也必随着增大。 \qty{100}{\celsius} 时，$K_w$ 的值是 \num{1e-12}，跟 \num{1e-14} 相比约增大 100 倍。
但在常温时，$K_w$ 的值一般可以认为是 \num{1e-14}。

\subsection{溶液的酸碱性和 pH 值}
$[\ce{H+}]$ 和 $[\ce{OH-}]$ 的乘积是一个常数 \num{1e-14}，由于水的电离平衡的存在，不仅纯水是这样，就是在酸性（或碱性）的稀溶液里也是这样。
在常温时，中性溶液里氢离子浓度和氢氧根离子浓度相等，都是 \qty{1e-7}{mol/L}; 在酸性溶液里不是没有 \ce{OH-}，只是含有的 \ce{H+} 多一些；在碱性溶液里也不是没有 \ce{H+}，只是含有的 \ce{OH-} 多一些。
总之，不管稀溶液是酸性、碱性或中性，在常温时，$[\ce{H+}]$ 和 $[\ce{OH-}]$ 的乘积都等于 \num{1e-14}。
由此可知，知道了溶液中 \ce{H+} 的浓度，就可以知道溶液中 \ce{OH-} 的浓度。
如已知常温时某溶液 \ce{H+} 浓度为 \qty{1e-1}{mol/L}，那么; 该溶液的 \ce{OH-} 浓度为：
\[ [\ce{OH-}] = \frac{K_w}{[\ce{H+}]}=\frac{\num{1e-14}}{\num{1e-1}}=\qty{e-13}{mol/L}\]

常温下，溶液的酸碱性跟 \ce{H+} 浓度和 \ce{OH-} 浓度的关系可以表示如下：

\begin{tblr}{colspec={ccl},hlines=0pt,vlines=0pt}
  中性溶液& & $[\ce{H+}]=[\ce{OH-}]=\qty{1e-7}{M}$ \\
  酸性溶液& & $[\ce{H+}]>[\ce{OH-}]$，$[\ce{H+}]>\qty{1e-7}{M}$ \\
  碱性溶液& & $[\ce{H+}]<[\ce{OH-}]$，$[\ce{H+}]<\qty{1e-7}{M}$ \\
\end{tblr}

\noindent
$[\ce{H+}]$ 越大，溶液的酸性越强; $[\ce{H+}]$ 越小，溶液的酸性越弱。
我们经常要用到一些 \ce{H+} 浓度很小的溶液，如 $[\ce{H+}]$ 等于  \qty{e-7}{M}、\qty{1.34e-3}{M} 等等。
用这样的数值来表示溶液酸碱性的强弱，很不方便。
为此化学上常采用 \ce{H+} 浓度的负对数来表示溶液酸碱性的强弱，叫做\Concept{溶液的 pH 值}。
\[ \text{pH} = -\lg[\ce{H+}]\]

例如：纯水的 $[\ce{H+}]=\qty{1e-7}{M}$，它的 pH 值是：
\[ \text{pH} = -\lg\num{e-7} =-(-7)=7 \]

又如：$[\ce{H+}]=\qty{e-5}{M}$ 的酸性溶液，它的 pH 值是：
\[ \text{pH} = -\lg\num{e-5} =5 \]

$[\ce{H+}]=\qty{e-9}{M}$ 的碱性溶液，它的 pH 值是：
\[ \text{pH} = -\lg\num{e-9} =9 \]

$[\ce{H+}]=\qty{1}{M}$ 的碱性溶液，它的 pH 值是：
\[ \text{pH} = -\lg[\ce{H+}]= -\lg 1 =0 \]

因此，在中性溶液里 pH 值等于 7；在酸性溶液里 pH 值小于 7；在碱性溶液里 pH 值大于 7。
溶液的酸性越强，pH 值越小。
溶液的碱性越强，pH 值越大。
所以，可以用 pH 值表示溶液的酸性强弱；同时，也可以用 pH 值表示溶液的碱性强弱，如\cref{fig:5-3} 所示。
\begin{figure}
  \includegraphics{5-3.pdf}
  \caption{$[\ce{H+}]$ 和 pH 值的对照关系的示意图}\label{fig:5-3}
\end{figure}

当溶液的 \ce{H+} 和 \ce{OH-} 的浓度大于 \qty{1}{M} 时，用 pH 值表示酸、碱性的强弱并不简便。如\cref{tab:5-3} 所示。
\begin{table}
  \caption{酸溶液里 \ce{H+} 浓度大于 \qty{1}{M} 时的 pH 值}\label{tab:5-3}
  \begin{tblr}{colspec={*4{X[c]}},hline{2}=0.8pt}
    \ce{[H+]} & $2M$ & $4M$ & $6M$ \\
    pH & $-0.3$ & $-0.6$ & $-0.78$ \\
  \end{tblr}
\end{table}

所以当溶液的 \ce{H+} 的浓度大于 \qty{1}{M} 时，一般不用 pH 值表示溶液的酸碱性，而是直接用 \ce{H+} 的浓度来表示。


\subsection{酸碱指示剂}
测定溶液 pH 值的方法很多，通常可用酸碱指示剂、pH 试纸或 pH 计（一种测定溶液 pH 值的仪器）。
酸碱指示剂一般是弱有机酸或弱有机碱。
它们的颜色变化是在一定的 pH 值范围内发生的。
我们把指示剂发生颜色变化的 pH 值范围叫做\Concept{指示剂的变色范围}。
各种指示剂的变色范围是由实验测定的。
从\cref{fig:5-4} 可以看到石蕊、酚酞和甲基橙的变色范围是各不相同的。
\begin{figure}
  \includegraphics{5-4.pdf}
  \caption{甲基橙、酚酞和石蕊的变色范围}\label{fig:5-4}
\end{figure}

pH 值的测定和控制在工农业生产、科学研究和医疗卫生中都很重要。测定溶液的 pH 值比较简便的方法是用 pH 试纸。
这种试纸是由多种指示剂的混和溶液浸制而成。
把待测试液滴在 pH 试纸上，试纸上显出的颜色跟标准比色卡相比，就可以知道该溶液的 pH 值。
测定溶液 pH 值最精确的方法是用 pH 计。

\begin{Theorem}{讨论}
  产生弱电解质的复分解反应为什么能够发生？试用电离平衡的观点来解释。
\end{Theorem}

\begin{Practice}[习题]
\begin{question}
  \item 计算在 \qty{25}{\celsius} 时，水的 $K_\text{电离}$ 是多少。
  \item \qty{0.2}{M} \ce{HCl} 和 $0.2M$ \ce{CH3COOH} 溶液中 $[\ce{H+}]$ 各是多少？
  \item 计算下列溶液中的 $[\ce{H+}]$ 和 $[\ce{OH-}]$：
  \begin{tasks}(2)
    \task $0.1M$ \ce{NaOH}，
    \task \qty{1e-3}{M} \ce{HCl}，
    \task $0.1M$ \ce{CH3COOH}，
    \task $0.1M$ \ce{NH3.H2O}。
  \end{tasks}
  \item 酸性水溶液里有没有 \ce{OH-}？碱性水溶液里有没有 \ce{H+}？为什么？
  \item 什么叫做 pH 值？水溶液的 pH 值跟溶液的酸碱性的关系怎样？
  \item 计算下列溶液的 pH 值：
  \begin{tasks}(2)
    \task $0.01M$ \ce{HCl}，
    \task $0.001M$ \ce{KOH}，
    \task $0.01M$ \ce{NH3.H2O}，
    \task $0.1M$ \ce{CH3COOH}，
    \task $[\ce{H+}]=\num{e-6}{M}$，
    \task $[\ce{OH-}]=\num{e-9}{M}$。
  \end{tasks}
  \item 已知某溶液的 $[\ce{H+}]$ 为 $\num{e-4}{M}$，计算该溶液的 $[\ce{OH-}]$ 和 pH 值。
  \item 有三种溶液 $A$、$B$、$C$，其中 $A$ 的 pH 值为 5，$B$ 中 $[\ce{H+}]=\num{e-4}{M}$，$C$ 中 $[\ce{OH-}]=\num{e-11}{M}$，问哪个酸性强？
  \item 在 \qty{1}{L} 溶液里含有 \ce{NaOH} \qty{4}{g}，求该溶液的 pH 值。
  \item 当某稀溶液的 pH 值增加 2 个单位时，$[\ce{H+}]$ 和 $[\ce{OH-}]$ 怎样改变？当 pH 值减少 3 个单位时，$[\ce{H+}]$ 和 $[\ce{OH-}]$ 又怎样改变？
  \item 健康人的血液的 pH 值为 \numrange{7.35}{7.45}，患某种疾病的人的血液 pH 值可暂时降到 5.9，问这时血液中氢离子浓度约为正常状态的多少倍？
  \item 什么叫做指示剂的变色范围？如某溶液的 pH 值是 6，分别滴入
  \begin{enumerate*} 
    \item 甲基橙、
    \item 石蕊、
    \item 酚酞试液，
  \end{enumerate*}溶液各显出什么颜色？
\end{question}
\end{Practice}

\section{盐类的水解}
\subsection{盐类的水解}
有的盐溶解于水后，所形成的水溶液能显出一定的酸碱性。
\begin{Experiment}\label{expr:5-2}
把少量 \ce{CH3COONa}、\ce{NH4Cl}、\ce{NaCl} 的晶体分别投入三个盛有蒸馏水的试管。振荡试管使之溶解。然后分别用 pH 试纸加以检验。
\end{Experiment}
实验结果表明: \ce{CH3COONa} 的水溶液显碱性，\ce{NH4Cl} 的水溶液显酸性，\ce{NaCl} 的水溶液显中性。这是什么原因呢？

我们学过水能微弱地电离出 \ce{H+} 和 \ce{OH-}，二者的浓度相等，并且处于动态平衡状态。

\cref{expr:5-2} 中所用的 \ce{CH3COONa} 是由一种强碱（氢氧化钠） 和一种弱酸（醋酸）中和所生成的盐。在它的水溶液里，并存着下列几种电离：
\[\ce{CH3COONa = CH3COO- + Na+}\tikz[remember picture,overlay]\coordinate(a);
\tikz[remember picture,overlay]{
  \node (b) at ([xshift=-6em,yshift=-4.8ex]a){\ce{H2O <=> H+ + OH-}};
  \node (p) at ([xshift=0.3em,yshift=0.6ex]b.north){$+$};
  \node (c) at ([yshift=-5.5ex]p.south)[rotate=90]{\ce{<=>}};
  \node (d) at (c.west)[]{\ce{CH3COOH}};}
\]
\par\vskip2.3cm
可以清楚地看到，由于 \ce{CH3COO-} 跟水里的 \ce{H+} 结合而生成难电离的 \ce{CH3COOH} 消耗了溶液中的 \ce{H+}，从而破坏了水的电离平衡。随着溶液里 \ce{H+} 的浓度减小，水的电离平衡向右移动，于是 \ce{OH-} 浓度随着增大，直至建立新的平衡。结果，溶液里 $[\ce{OH-}]>[\ce{H+}]$，从而使溶液显出碱性。上述反应可用离子方程式表示如下：
\[ \ce{ CH3COOH- + H2O <=> CH3COOH + OH- }\]

这种\emph{在溶液中盐的离子跟水所电离出来的 \ce{H+} 或 \ce{OH-} 生成弱电解质的反应，叫做}\Concept{盐类的水解}。

从上式可见，盐类水解后生成了酸和碱，即盐类的水解反应可看作是酸碱中和反应的逆反应：
\[ \ce{ \text{酸} + \text{碱} <=>[\text{中和}][\text{水解}] \text{盐} + \text{水} } \]

盐类的水解跟生成这种盐的酸和碱的强弱有着密切的关系。现在分别说明如下：
\begin{enumerate}
  \item 强碱和弱酸所生成盐的水解
  
  上面所讨论的醋酸钠就是由强碱（氢氧化钠）和弱酸（醋酸）所生成的盐，这种盐水解后使溶液显碱性。

碳酸钠也是由强碱（氢氧化钠）和弱酸（碳酸）所生成的盐，它水解后，溶液也显碱性。由于碳酸是二元酸，所以碳酸钠的水解反应复杂一些，要分两步进行:

第一步是碳酸钠在水溶液里电离出来的 \ce{CO3^2-} 进行水解。
\[\ce{Na2CO3 = 2Na+ + CO3^2-}\tikz[remember picture,overlay]\coordinate(a);
\tikz[remember picture,overlay]{
  \node (b) at ([xshift=-5.7em,yshift=-4.8ex]a){\ce{H2O <=> OH- + H+}};
  \node (p) at ([xshift=-1.4em,yshift=0.6ex]b.north east){$+$};
  \node (c) at ([yshift=-5.5ex]p.south)[rotate=90]{\ce{<=>}};
  \node (d) at (c.west)[]{\ce{HCO3-}};}
\]
\par\vskip2.3cm
离子方程式是：$\ce{CO3^2- + H2O <=> HCO3- + OH-}$

第二步是生成的 \ce{HCO3-} 进一步进行水解。

离子方程式是：$\ce{HCO3^- + H2O <=> H2CO3 + OH-}$

由此可见，溶液里的 \ce{CO3^2-} 跟由水分子电离出来的 \ce{H+} 结合生成 \ce{HCO3-}，\ce{HCO3-} 又跟 \ce{H+} 结合生成 \ce{H2CO3}，促使水继续电离，溶液里的 \ce{OH-} 浓度增加，所以溶液显碱性。

但是，\ce{Na2CO3} 第二步水解的程度很小，平衡时溶液中 \ce{H2CO3} 分子的浓度很小，不会放出 \ce{CO2} 气体。

其它如碳酸钾、硫化钠、磷酸钠等盐的水解也属于这种类型。
  \item 强酸和弱碱所生成盐的水解
  
  \cref{expr:5-2} 中所用的 \ce{NH4Cl} 就是由强酸（盐酸）和弱碱 （氨水）中和所生成的盐。它在水溶液里的水解过程表示如下：
  \[\ce{NH4Cl = NH4+ + Cl-}\tikz[remember picture,overlay]\coordinate(a);
  \tikz[remember picture,overlay]{
    \node (b) at ([xshift=-5em,yshift=-4.8ex]a){\ce{H2O <=> OH- + H+}};
    \node (p) at ([xshift=0.5em,yshift=0.6ex]b.north){$+$};
    \node (c) at ([yshift=-5.5ex]p.south)[rotate=90]{\ce{<=>}};
    \node (d) at (c.west){\ce{NH3.H2O}};}
  \]
  \par\vskip2.3cm
  在这里，由于 \ce{NH4+} 跟水里的 \ce{OH-} 结合而生成难电离的 \ce{NH3.H2O}，打破了水的电离平衡。随着溶液里 \ce{OH-} 浓度减小，水的电离平衡也将向右移动，于是 \ce{H+} 浓度随着增大，直至建立新的平衡。结果，溶液里的 \ce{H+} 浓度大于 \ce{OH-} 浓度，从而使溶液显出酸性。这一反应也可以用离子方程式来表示：
  \[ \ce{ NH4+ + H2O <=> NH3.H2O + H+ }\]

  其它如硝酸铜、硫酸铜、硫酸铵等盐的水解都属于这种类型。
  \item 弱酸和弱碱所生成盐的水解
  
  由弱酸和弱碱所生成的盐在水溶液里也会起水解反应。例如 \ce{CH3COONH4} 是弱酸（醋酸）和弱碱（氨水）所生成的盐，它在水里的水解过程表示如下：
  \[\ce{CH3COONH4 = CH3COO- + NH4+ }\tikz[remember picture,overlay]\coordinate(a);
  \tikz[remember picture,overlay]{
    \node (b)  at ([xshift=-6em,yshift=-4.8ex]a){\ce{H2O <=> H+} \quad $+$ \quad \ce{OH-}};
    \node (p1) at ([xshift=-0.9em,yshift=0.6ex]b.north){$+$};
    \node (c1) at ([yshift=-5.5ex]p1.south)[rotate=90]{\ce{<=>}};
    \node (d1) at (c1.west){\ce{CH3COOH}};
    \node (p2) at ([xshift=4.5em,yshift=0.6ex]b.north){$+$};
    \node (c2) at ([yshift=-5.5ex]p2.south)[rotate=90]{\ce{<=>}};
    \node (d2) at (c2.west){\ce{NH3.H2O}};
  }
  \]
  \par\vskip2.3cm

  上述反应可用离子方程式表示如下：
  \[ \ce{ CH3COO- + NH4+ + H2O <=> CH3COOH + NH3.H2O } \]

  很清楚，由于这种盐的弱酸根和弱碱根能分别跟水里的 \ce{OH-} 和 \ce{OH-} 结合而生成难电离的 \ce{CH3COOH} 和 \ce{NH3.H2O}，破坏了水的电离平衡，从而使水的电离向右移动。
至于这类盐的水溶液显酸性还是碱性，要取决于所生成的醋酸和氨水的电离常数的相对大小。
如果是弱酸的电离常数大，那么溶液显酸性；如果是弱碱的电离常数大，那么溶液显碱性；如果两者的电离常数相等，那么溶液呈中性。
例如 \ce{CH3COOH} 和 \ce{NH3.H2O} 的电离常数分别是 \num{1.75e-5} 和 \num{1.77e-5}，基本相等，所以 \ce{CH3COONH4} 的水溶液接近中性。

又如 \ce{NH4CN} 是弱碱（氨水）跟弱酸（氢氰酸）所生成的盐，由于氨水的电离常数（\qty{25}{\celsius} 时为 \num{1.75e-5}）大于氢氰酸的电离常数（\qty{25}{\celsius} 时为 \num{4.93e-10}），所以当 \ce{NH4CN} 水解时溶液显碱性。
\end{enumerate}

上述三种类型的盐能够发生水解，基本原因在于组成盐的离子能跟水电离出来的 \ce{H+} 或 \ce{OH-} 结合形成了弱电解质。

强酸和强碱所生成的盐，如 \ce{NaCl} 等，因为它们电离生成的阴、阳离子，都不跟溶液中的 \ce{H+} 或 \ce{OH-} 结合形成弱电解质，所以水中的 \ce{H+} 和 \ce{OH-} 的数目保持不变，没有破坏水的电离平衡。
因此，这种由强酸和强碱所生成的盐不发生水解，溶液显中性。
\cref{expr:5-2} 中所用的 \ce{NaCl}，就是这种盐的例子。
其它如氯化钾、硫酸钠、硝酸钠等盐不发生水解都是由于这个原因。

\subsection{盐类水解的利用}
盐类水解程度的大小，主要由盐的本性所决定。
同时当水解达到平衡时，跟其它平衡一样，也受温度、浓度等的影响。

盐类水解是中和反应的逆反应。
中和反应是放热反应，所以水解必然是吸热反应。
因此，升高温度能促进盐类的水解。
如用纯碱水溶液洗涤油污物品时，热的碱水去油污的效果就比较好。
这是因为加热能促进纯碱水解，使 \ce{OH-} 浓度增大，所以增强了去油污能力。

当增大或减小 \ce{H+} 或 \ce{OH-} 的浓度时，平衡要向左或向右移动，这样可以抑制或促进水解反应的进行。
如在实验室配制 \ce{FeCl3} 溶液时，由于 \ce{FeCl3} 是强酸弱碱生成的盐，容易水解生成难溶于水的 \ce{Fe(OH)3}：
\[ \ce{ FeCl3 + 3H2O <=> Fe(OH)3 + 3HCl }\]
致使溶液混浊，得不到澄清的 \ce{FeCl3} 溶液。
所以配制 \ce{FeCl3} 溶液时，为了防止水解，通常要向溶液中加入一定量的盐酸。

\begin{Theorem}{讨论}
  盐是中和反应的产物，为什么某些盐又会发生水解反应（中和反应的逆反应）？试用电离平衡的观点加以解释。
\end{Theorem}

\begin{Practice}[习题]
\begin{question}
  \item 下列几种盐，哪些能水解，哪些不能水解？它们溶液的酸、碱性怎样？写出水解反应的离子方程式。
  \begin{tasks}(4)
    \task \ce{NH4NO3}；
    \task \ce{CH3COOK}；
    \task \ce{NaNO3}；
    \task \ce{FeSO4}。
  \end{tasks}
  \item 怎样用最简单的方法区别下列三种溶液: \ce{NaCl} 溶液、 \ce{NH4Cl} 溶液和 \ce{Na2CO3} 溶液。
  \item 如果要使 \ce{FeCl3} 的水解能进行得较为完全，可以采取哪些措施？
  \item 草木灰是农村常用的钾肥，它的主要成分是碳酸钾。解释为什么草木灰不宜与用作氮肥的铵盐混和使用。
  \item 泡沫灭火器中盛有 \ce{Al2(SO4)3} 和 \ce{NaHCO3} 溶液，为什么把它们混和后就能有灭火作用？这里发生了哪些化学反应？
  \item 食盐溶液和醋酸铵溶液都呈中性，但它们呈中性的原因有着本质的区别。试加以说明。
\end{question}
\end{Practice}

\section{酸碱的当量浓度}
\subsection{酸和碱的克当量}
我们在初中化学里已经学过酸和碱的溶液在一起能发生中和反应，现在进一步来研究它们在反应中的量的关系。下列反应都是中和反应：
\begin{gather*}
  \ce{ NaOH + HCl \xlongequal{\quad} NaCl + H2O }\\
  \ce{ 2NaOH + H2SO4 \xlongequal{\quad} Na2SO4 + 2H2O }
\end{gather*}
或
\begin{gather*}
  \ce{ NaOH + \frac{1}{2} H2SO4 \xlongequal{\quad} \frac{1}{2} Na2SO4 + H2O }\\
  \ce{ 3NaOH + H3PO4 \xlongequal{\quad} Na3PO4 + 3H2O }
\end{gather*}
或
\[ \ce{ NaOH + \frac13 H3PO4 \xlongequal{\quad} \frac13 Na3PO4 + H2O  }\]

从上面的化学方程式可以看出：一元酸、二元酸或三元酸跟 \qty{1}{mol} \ce{NaOH} 完全中和所需要酸的摩尔数是不相同的。
一元酸（如 \ce{HCl}）要用 \qty{1}{mol}，二元酸（如 \ce{H2SO4}）只要用 $\frac{1}{2}\,\unit{mol}$，三元酸（如 \ce{H3PO4}）只要用 $\frac{1}{3}\,\unit{mol}$。这是由于酸和碱的中和反应的实质就是：
\[ \ce{ H+ + OH- <=> H2O }\]
即 \qty{1}{mol} \ce{H+} 恰好能跟 \qty{1}{mol} \ce{OH-} 中和生成水。
而相同摩尔数的一元酸、二元酸或三元酸，在中和反应时所提供的 \ce{H+} 的摩尔数是不同的。
例如，\qty{1}{mol} \ce{HCl} 能提供 \qty{1}{mol} \ce{H+}，\qty{1}{mol} \ce{H2SO4} 能提供 \qty{2}{mol} \ce{H+}，\qty{1}{mol} \ce{H3PO4} 能提供 \qty{3}{mol} \ce{H+}。
为此，要完全中和 \qty{1}{mol} \ce{OH-}，如果用一元酸就需要 \qty{1}{mol}，如果用二元酸就只需要 $\frac12\,\unit{mol}$，用三元酸就只需要 $\frac13\,\unit{mol}$。
从具体的质量来说，\qty{1}{mol} 的盐酸的质量是 \qty{36.5}{g}，$\frac12\,\unit{mol}$ \ce{H2SO4} 的质量是 \qty{49}{g}，$\frac13\,\unit{mol}$ \ce{H3PO4} 的质量是 \qty{32.7}{g}。 
它们的质量虽然不等，但却具有相同的中和能力，都能跟 \qty{40}{g} \ce{NaOH} 完全中和。
化学上把这些反应中彼此相当的酸和碱的质量（\unit{g}），叫做酸和碱的克当量。
上述 \ce{HCl}、\ce{H2SO4}、\ce{H3PO4} 和 \ce{NaOH} 的克当量分别是 \qty{36.5}{g}、\qty{49}{g}、\qty{32.7}{g} 和 \qty{40}{g}。

酸和碱的克当量可用公式表示如下：
\begin{align*}
  \text{酸的克当量}&=\frac{\qty{1}{mol} \text{ 酸的质量}(\unit{g})}{\qty{1}{mol}\text{ 酸所提供}\ce{H+}\text{ 的摩尔数}}\\[10pt]
  \text{碱的克当量}&=\frac{\qty{1}{mol} \text{ 碱的质量}(\unit{g})}{\qty{1}{mol}\text{ 碱所提供}\ce{OH-}\text{ 的摩尔数}}\\
\end{align*}

例如，\qty{1}{mol} \ce{HNO3} 的质量是 \qty{63}{g}，\qty{1}{mol} \ce{HNO3} 能提供 \qty{1}{mol} \ce{H+}。
\[ \ce{HNO3} \text{ 的克当量}=\frac{63}{1}=\qty{63}{g}\]

又如，\qty{1}{mol} \ce{H2SO4} 的质量是 \qty{98}{g}，\qty{1}{mol} \ce{H2SO4} 能提供 \qty{2}{mol} \ce{H+}。
\[ \ce{H2SO4} \text{ 的克当量}=\frac{98}{2}=\qty{49}{g}\]

同理: \qty{1}{mol} \ce{CaCO3} 的质量是 \qty{74}{g}，\qty{1}{mol} \ce{Ca(OH)2} 能提供 \qty{2}{mol} \ce{OH-}。
\[ \ce{Ca(OH)2} \text{ 的克当量}=\frac{74}{2}=\qty{37}{g}\]

要注意在不同的化学反应中，酸或碱可以具有不同的克当量。例如：
\begin{gather*}
  \ce{ H2SO4 +  NaCl \xlongequal{\text{微热}} NaHSO4 +  HCl ^ }\\
  \ce{ H2SO4 + 2NaCl \xlongequal{\text{强热}} Na2SO4 + 2HCl ^ }
\end{gather*}

在第一个反应里，\qty{1}{mol} \ce{H2SO4} 只提供 \qty{1}{mol} \ce{H+}，所以 \ce{H2SO4} 的克当量为 \qty{98}{g}；在第二个反应里，\qty{1}{mol} \ce{H2SO4} 提供 \qty{2}{mol} \ce{H+}，所以 \ce{H2SO4} 的克当量是 $98/2=\qty{49}{g}$。
因此多元酸的克当量要根据具体反应来确定。
多元碱的克当量也要根据具体的反应来确定。

上面已经谈到 \qty{1}{mol} \ce{NaOH} 跟一元酸、二元酸、三元酸完全作用时，酸的摩尔数不一定相等，但 1 克当量 \ce{NaOH} 跟 \ce{HCl}、\ce{H2SO4}、\ce{H3PO4} 完全反应的量都是 1 克当量，即酸碱的克当量数一定相等。
酸或碱的克当量数、克当量和质量的关系，可用公式表示如下。
\[ \text{克当量数}=\frac{\text{酸或（碱）的质量 }\,(\unit{g})}{\text{酸或（碱）的克当量 }\,(\unit{g})}\]

\begin{example}
  计算 \qty{147}{g} \ce{H2SO4} 的克当量数。
\end{example}
\begin{solution}
  \ce{H2SO4} 的克当量是 \qty{49}{g}。
  \[ \ce{H2SO4}\text{ 的克当量数}=\frac{147}{49}=3\]
  答：\qty{147}{g} \ce{H2SO4} 的克当量数是 3。
\end{solution}

\begin{example}
  计算 \qty{20}{g} \ce{NaOH} 的克当量数。
\end{example}
\begin{solution}
  \ce{NaOH} 的克当量是 \qty{40}{g}。
  \[ \ce{NaOH}\text{ 的克当量数}=\frac{20}{40}=0.5\]
  答：\qty{20}{g} \ce{NaOH} 的克当量数是 0.5。
\end{solution}

\subsection{当量浓度}
当量浓度是表示溶液浓度的又一种方法。
它是用 \qty{1}{L}（\qty{1000}{ml}）溶液中所含溶质的克当量数来表示的溶液浓度，常用 $N$ 表示。
例如，在 \qty{1}{L} 溶液里含有 2 克当量溶质时，这个溶液的浓度就是 2 当量浓度，即用 $2N$ 表示。
溶液的当量浓度跟溶质的克当量数和溶液体积的关系可以用下式表示：
\[ \text{溶液的当量浓度 }(N)=\frac{\text{溶质的克当量数}}{\text{溶液的体积 }\,(\unit{L})}\]
或
\[ \text{溶质的克当量数} = \text{溶液的当量浓度 }(N)\times\text{溶液的体积 }\,(\unit{L}) \]

我们已经知道，在中和反应中，酸碱完全作用时，它们的克当量数一定相等。即：
\[ \text{酸的克当量数}=\text{碱的克当量数} \]
如果用 $N_1$、$N_2$ 分别表示酸碱两种溶液的当量浓度。用 $V_1$、$V_2$ 分别表示酸、碱两种溶液的体积（\unit{L}），可以得出下面的公式：
\begin{gather*}
  N_1V_1=N_2V_2\\
  \text{当} \quad N_1=N_2 \text{ 时，} V_1=V_2 \\
  \text{当} \quad N_1>N_2 \text{ 时，} V_1<V_2 
\end{gather*}

注意，这个公式里的 $V_1$ 和 $V_2$ 的单位不一定用升，只要二者相同，用其它体积单位也可以。

\subsection{有关当量浓度的计算}
\subsubsection{配制当量浓度溶液的计算}
\begin{example}
要配制 $0.1N$ \ce{NaOH} 溶液 \qty{100}{ml}，需用 \ce{NaOH} 多少克？
\end{example}

\begin{analyze}
  根据 \qty{100}{ml} $0.1N$ \ce{NaOH} 溶液的克当量数，就可以算出所需 \ce{NaOH} 的克数。
\end{analyze}

\begin{solution}
已知 \ce{NaOH} 的当量浓度是 $0.1N$，\ce{NaOH} 溶液的体积为 \qty{100}{ml}。

设 $x$ 为 \qty{100}{ml} $0.1N$ 的 \ce{NaOH} 溶液中含有 \ce{NaOH} 的克当量数。
\begin{gather*} 
  x=\frac{100\times 0.1}{1000}=0.01\\
  \ce{NaOH}\text{ 的质量}=\text{克当量数}\times\text{克当量}=0.01\times 40=\qty{0.4}{g} 
\end{gather*}
答: 需用 \qty{0.4}{g} \ce{NaOH} 才能配制 \qty{100}{ml} 的 $0.1N$ \ce{NaOH} 溶液。
\end{solution}

\subsubsection{酸碱中和反应的有关计算}
\begin{example}
有未知浓度的 \ce{H2SO4} 溶液 \qty{40}{ml}，需加入 $0.4N$ 的 \ce{NaOH} 溶液 \qty{24}{ml} 才能完全中和。问 \ce{H2SO4} 溶液的当量浓度是多少？ 这 \qty{40}{ml} 的 \ce{H2SO4} 溶液里含有 \ce{H2SO4} 多少克？
\end{example}

\begin{analyze}
先计算溶液中 \ce{NaOH} 的克当量数，然后运用酸碱溶液中和时克当量数相等的规律，即 $N_1V_1=N_2V_2$ 的公式，求得 \ce{H2SO4} 的当量浓度。
再由 \ce{H2SO4} 的克当量数算出 \qty{40}{ml} \ce{H2SO4} 溶液中所含 \ce{H2SO4} 的质量。
\end{analyze}

\begin{solution}
  已知 $N_1=0.4N$，$V_1=\qty{24}{ml}$，$V_2=\qty{40}{ml}$。

  设 $N_2$ 为 \ce{H2SO4} 的当量浓度。
  \begin{align*} 
    N_1V_1&=N_2V_2\\ 
    N_2&=\frac{N_1V_1}{V_2}=\frac{0.4\times24}{40}=0.24N
  \end{align*}

  \qty{40}{ml} $0.24N$ \ce{H2SO4} 溶液中含纯 \ce{H2SO4} 的质量
  \[ 49\times\frac{40}{1000}\times 24=\qty{0.47}{g} \]
  答: \ce{H2SO4} 溶液的当量浓度是 $0.24N$，\qty{40}{ml} \ce{H2SO4} 溶液中含纯 \ce{H2SO4} \qty{0.47}{g}。
\end{solution}

\subsubsection{计算溶液在稀释前后的体积和当量浓度}
可用 $V_1$、$N_1$ 和 $V_2$、$N_2$ 分别代表稀释前后溶液的体积和当量浓度。
然后利用稀释前后溶质克当量数相等的公式 $N_1V_1=N_2V_2$ 即可算出溶液在稀释前后的体积和当量浓度。

\begin{example}
  用 $12N$ 的浓盐酸配制成 $6N$ 的稀盐酸 \qty{500}{ml}，要用 $12N$ 的浓盐酸多少毫升？
\end{example}
\begin{solution}
  已知 $N_1=6N$，$V_1=\qty{500}{ml}$，$N_2=12N$。

  设 $V_2$ 为要用 $12N$ 的浓盐酸的毫升数。
  \begin{align*} 
    12V_2&=6\times 500\\ 
    V_2&=\qty{250}{ml}
  \end{align*}
  答: 要用 $12N$ 的浓盐酸 \qty{250}{ml}。
\end{solution}

\begin{example}
  用 $12N$ 的浓盐酸 \qty{100}{ml} 配制成 \qty{600}{ml} 稀盐酸的当量浓度是多少？
\end{example}
\begin{solution}
  已知 $N_1=12N$，$V_1=\qty{100}{ml}$，$V_2=\qty{600}{ml}$。

  设 $N$ 为稀盐酸的当量浓度
  \begin{align*} 
    12\times100&=600N_2\\ 
    N_2&=2N
  \end{align*}
  答: 配制成的稀盐酸的当量浓度是 $2N$。
\end{solution}

\subsubsection{当量浓度和百分比浓度的换算}
\begin{example}
  现有 98\% 的 \ce{H2SO4} 的密度是 \qty{1.84}{g/ml}，要配制 \qty{500}{ml} $0.50N$ 的稀 \ce{H2SO4} 溶液，需要这样的 \ce{H2SO4} 多少毫升？
\end{example}

\begin{analyze}
  先把 98\% 的 \ce{H2SO4} 换算为当量浓度，然后再用 $N_1V_1=N_2V_2$ 这个公式计算这种 \ce{H2SO4} 的体积。
\end{analyze}

\begin{solution}
  已知 $V_1=\qty{500}{ml}$，$N_1=0.50N$。
  \[ N_2=\frac{1.84\times1000\times 98\%}{49}=36.8 N\]

  设 $V_2$ 为 98\% \ce{H2SO4} 的毫升数。
  \begin{align*} 
    36.8\times V_2&=0.5\times 500\\ 
    V_2&=\qty{6.8}{ml}
  \end{align*}
  答: 配制 \qty{500}{ml} $0.5N$ 稀硫酸需要 98\% 的浓硫酸 \qty{6.8}{ml}。
\end{solution}

\begin{Practice}[习题]
\begin{question}
  \item 求下列化合物的充当量：
  \begin{tasks}(3)
    \task \ce{CH3COOH}
    \task \ce{KOH}
    \task \ce{Mg(OH)2}
  \end{tasks}
  \item 求 \qty{1}{mol} 的下列化合物所含的克当量数：
  \begin{tasks}(3)
    \task \ce{H2SO4}
    \task \ce{H3PO4}
    \task \ce{HClO4}
  \end{tasks}
  \item 下列 \qty{50}{ml} $0.4N$ 的各溶液中，含有溶质各多少克？
  \begin{tasks}(3)
    \task \ce{H2SO4} 溶液
    \task \ce{NaOH} 溶液
    \task \ce{HCl} 溶液
  \end{tasks}
  \item 计算下列常用试剂的当量浓度：
  \begin{tasks}[label=(\arabic*)]
    \task 盐酸——密度 \qty{1.19}{g/cm^3}，含 \ce{HCl} 37\%，
    \task 硝酸——密度 \qty{1.42}{g/cm^3}，含 \ce{HNO3} 71\%，
    \task 硫酸——密度 \qty{1.84}{g/cm^3}，含 \ce{H2SO4} 98\%。
  \end{tasks}
  \item $2N$ \ce{NaOH} 溶液的密度是 \qty{1.08}{g/cm^3}，计算出溶液的百分比浓度。
  \item 50\% 的 \ce{NaOH} 溶液的密度为 \qty{1.525}{g/cm^3}，计算该溶液的当量浓度。
  \item 配制 $0.2N$ \ce{HCl} 溶液 \qty{1.5}{L}，需要 $6N$ \ce{HCl} 多少毫升？
  \item 要中和 \qty{25}{ml} $0.1N$ \ce{NaOH} 溶液，需要 $0.2N$ \ce{H2SO4} 溶液多少毫升？
  \item $0.5N$ \ce{NaOH} 溶液 \qty{20}{ml}，恰好跟 \qty{10}{ml} \ce{HNO3} 溶液中和，求这种 \ce{HNO3} 溶液的当量浓度和所含 \ce{HNO3} 的质量。
  \item 下列两种情形是否可能？为什么？举例说明。
  \begin{tasks}[label=(\arabic*)]
    \task 体积相同、摩尔浓度不同的两种酸溶液，能否用同体积、同摩尔浓度的同一种碱溶液使它们完全中和。
    \task 体积相同、当量浓度不同的两种酸溶液，能否用同体积、同当量浓度的同一种碱溶液使它们完全中和。
  \end{tasks}
  \item 把 $5N$ \ce{NaOH} 溶液 \qty{200}{ml} 稀释到 \qty{1}{L}，它的当量浓度是多少？
  \item 现有 \qty{100}{ml} $0.50N$ 的 \ce{NaOH} 溶液，要加入多少克水才能使溶液的当量浓度降低为 $0.10N$？
  \item 在 \qty{100}{ml} $1N$ \ce{H2SO4} 溶液中投入锌粒，锌完全反应后，还须用 $0.5N$ \ce{NaOH} 溶液 \qty{30}{ml} 才能完全中和剩余的硫酸。求锌粒的质量是多少克？
\end{question}
\end{Practice}

\section{酸和碱的中和反应}
\subsection{酸碱中和滴定}
我们在初中化学里已经学过怎样在氢氧化钠溶液里滴加盐酸，使我们初步懂得了使酸和碱发生反应达到中和的方法。现在我们要在学习酸碱当量浓度知识的基础上，来进一步掌握用已知当量浓度的酸（或碱），来测定未知当量浓度的碱（或酸）的方法，这就是\Concept{酸碱中和滴定}。

现在让我们用已知当量浓度的 \ce{HCl} 溶液来滴定未知当量浓度的 \ce{NaOH} 溶液为例，以说明酸碱中和滴定的过程。
\par\medskip\noindent
\begin{minipage}{0.70\linewidth}\parindent2em
把 $0.1N$ 的 \ce{HCl} 溶液，注入洁净的酸式滴定管到刻度 “0” 以上。
把滴定管固定在滴定管夹上，如\cref{fig:5-5} 所示。
轻轻转动下面的活塞，把管的尖嘴部分充满溶液，然后调整管内液面，使其保持在 “0” 或 “0” 以下的某一定的刻度，并记下准确读数; 把待测浓度的氢氧化钠溶液注入碱式滴定管，也把它固定在滴定管夹上。
轻轻挤压玻璃球，使管的尖嘴部分充满溶液，然后调整管内液面，使其保持在 “0” 或 “0” 以下某一定刻度，并记下准确读数。

然后在管下放一洁净的锥形瓶。
从碱式滴定管放出 \qty{20}{ml} \ce{NaOH} 溶液，注入锥形瓶，加入 2 滴酚酞指示剂，溶液立即呈红色。
然后，把锥形瓶移在酸式滴定管下，按\cref{fig:5-5} 所示进行操作，逐滴加入 $0.1N$ 的 \ce{HCl} 溶液，同时不断摇动锥形瓶，使溶液充分混和。
随着 \ce{HCl} 溶液的逐滴加入，锥形瓶里碱溶液的 \ce{OH-} 的浓度逐渐减少，因而溶液的红色逐渐消褪。
最后，当我们看到加入一滴 \ce{HCl} 溶液时，溶液立即褪成无色; 再加入一滴 \ce{NaOH} 溶液时，溶液又显红色，这就说明溶液里酸和碱的克当量数已经相等。
这时应该立即停止滴定。
然后准确记下两个滴定管上的刻度，并记下准确读数，求得用去的 \ce{NaOH} 和 \ce{HCl} 溶液的体积。
然后根据 $N_1V_1=N_2V_2$ 的公式，算出溶液中碱的含量。
\end{minipage}\hfill
\begin{minipage}{0.28\linewidth}
\begin{figurehere}
  \includegraphics{5-5.pdf}
  \caption{中和滴定的装置和操作示意图}\label{fig:5-5}
\end{figurehere}
\end{minipage}\par\medskip

\begin{Reading}[补充阅读]{}
在中和滴定过程中，酸和碱完全起反应时，即酸和碱的克当量数恰好相等时，化学上称它作等当点。
在等当点时溶液不一定都是中性的，有时显酸性，有时显碱性，这要看中和后生成的盐的性质来决定。
因为有的盐会起水解反应的缘故。
上面滴定用的是强酸强碱，所生成的盐不水解，所以在等当点时溶液是中性的。
\end{Reading}

% \begin{figure}
%   \caption{中和滴定的装置和操作示意图}\label{fig:5-5}
% \end{figure}

\subsection{中和热}
我们知道，化学反应都伴随着能量的变化，并且通常表现为热量的变化，即有吸热或放热现象发生。
实验测出，酸、碱在发生中和反应时有热量放出。
当 \qty{1}{L} $1M$ 的盐酸跟 \qty{1}{L} $1M$ 的氢氧化钠溶液起中和反应时，能放出 \qty{13.7}{kCal} 的热量。
\[ \ce{ NaOH \text{（稀）} + HCl \text{（稀）} \xlongequal{\quad} NaCl \text{（稀）} + H2O } + \qty{13.7}{kCal}\]

如果用 \qty{1}{L} $1M$ 的氢氧化钾溶液中和 \qty{1}{L} $1M$ 硝酸溶液，也放出 \qty{13.7}{kCal} 的热量。
\[ \ce{ KOH \text{（稀）} + HNO3 \text{（稀）} \xlongequal{\quad} KNO3 \text{（稀）} + H2O } + \qty{13.7}{kCal}\]

\emph{在稀溶液中，酸跟碱发生中和反应而生成 1 摩尔水，这时的反应热就是}\Concept{中和热}。

我们早已知道，中和反应的实质是 \ce{H+} 跟 \ce{OH-} 化合生成水，因此，上述反应可以用同一的离子方程式表示。
\[ \ce{ H+ + OH- \xlongequal{\quad} H2O } \]

当强酸跟强碱在稀溶液中发生中和反应，\qty{1}{mol} \ce{H+} 跟 \qty{1}{mol} \ce{OH-} 起反应而生成 \qty{1}{mol} 水，都放出 \qty{13.7}{kCal} 的热量。

\begin{Practice}[习题]
\begin{question}
  \item 为什么用中和滴定的方法可以测定未知当量浓度的酸或碱？
  \item 在酸碱中和滴定操作中，指示剂起什么作用？
  \item 为什么在进行滴定操作之前，一定要使溶液充满滴定管的尖嘴部分？
  \item 计算下列中和反应里放出的热量。
  \begin{tasks}[label=(\arabic*)]
    \task 用 \qty{20}{g} 氢氧化钠配成的稀溶液跟足量的稀盐酸起反应，能放出多少千卡的热量？
    \task 用 \qty{28}{g} 氢氧化钾配成的稀溶液跟足量的稀硝酸起反应，能放出多少千卡的热量？
  \end{tasks}
  \item 用稀 \ce{HCl} 和稀 \ce{HNO3} 的混和溶液 \qty{20}{ml} 滴入 \qty{10.3}{ml} $1N$ \ce{NaOH} 溶液恰好中和。如往这种混和溶液里加入过量的 \ce{AgNO3}，所生成的 \ce{AgCl} 沉淀经洗涤、烘干后重 \qty{0.98}{g}。求 \qty{1}{L} 这种混和溶液里含纯 \ce{HCl} 和纯 \ce{HNO3} 各多少克。
\end{question}
\end{Practice}


\section{原电池\texorpdfstring{\quad}{ }金属的腐蚀和防护}\label{sec:battery}
\subsection{原电池}
我们知道，物质发生化学反应时常伴有化学能跟热能、光能等的相互转化。
如在一般化学反应里，常表现出放热或吸热，有的化学反应还表现出发光，等等。
现在我们要来研究化学能是怎样转变为电能的。
\begin{Experiment}*[righthand ratio=0.5]
  把一块锌片和一块铜片平行地插入盛有稀硫酸溶液的烧杯里，可以看到锌片上有气体放出，铜片上没有气体放出。再用导线把锌片和铜片连接起来（\cref{fig:5-6}），观察铜片上有没有气体放出？在导线中间接入一个电流计，观察指针是否偏转？
  \tcblower
  \begin{figurehere}
    \includegraphics{5-6.pdf}
    \caption{原电池示意图}\label{fig:5-6}
  \end{figurehere}
\end{Experiment}
实验结果表明，用导线连接后，锌片在不断溶解，铜片上有氢气产生。
电流计上指针发生偏转。
这说明当铜片和锌片一同浸入稀 \ce{H2SO4} 时，由于锌比铜活泼，容易失去电子，锌被氧化成 \ce{Zn^2+} 而进入溶液，电子由锌片通过导线流向铜片，溶液中的 \ce{H+} 从铜片获得电子，被还原成氢原子，氢原子结合成氢分子从铜片上放出。变化过程可以表示如下：
\begin{gather*}
  \text{锌片}\quad \ce{ Zn - 2e = Zn^2+ } \quad\text{（还原反应）}\\
  \text{铜片}\quad \ce{ 2H+ +2e = 2H } \quad\text{（还原反应）}\\
  \ce{ 2H = H2 ^ }
\end{gather*}

这个实验充分证明，上述氧化—还原反应确实因电子的转移而产生电流。这种\emph{把化学能转变为电能的装置叫做}\Concept{原电池}（\cref{fig:5-6}）。
原电池中电子流出的一极是负极（如锌片），电极被氧化。
电子流入的一极是正极（如铜片），\ce{H+} 在电极上被还原。

人们应用原电池的原理制作了多种电池，如干电池、蓄电池以及供人造地球卫星、宇宙火箭、空间电视转播站使用的高能电池，等等。
电池在生活、工农业生产以及科学技术等方面都有广泛的用途。

但是，原电池的反应也使金属腐蚀而被损耗。
下面我们来学习金属的腐蚀和防护知识。

\subsection{金属的腐蚀和防护}
金属腐蚀是指金属或合金跟周围接触到的气体或液体进行化学反应而腐蚀损耗的过程。
比较普遍的是不纯的金属（或合金），接触到电解质溶液发生原电池反应，比较活泼的金属原子失去电子而被氧化所引起的腐蚀，这种腐蚀叫做\Concept{电化腐蚀}。
钢铁在潮湿的空气里所发生的腐蚀，就是电化腐蚀的最普通的例子。

我们都知道，钢铁在干燥的空气里长时间不易腐蚀，但在潮湿的空气里却很快地就会腐蚀，这是什么原因呢？
原来在潮湿的空气里，钢铁表面吸附一层薄薄的水膜而促使钢铁腐蚀。
水是弱电解质，它能电离出少量的 \ce{H+} 和 \ce{OH-}，同时由于空气里的二氧化碳的溶解，使水里的 \ce{H+} 增多。
\[ \ce{ H2O + CO2 <=> H2CO3 <=> H+ + HCO3- } \]
结果在钢铁表面形成了一层电解质溶液的薄膜，它跟钢铁里的铁和少量的碳恰好构成了原电池。
因此，这些钢铁制品的表面就形成了无数微小的原电池（\cref{fig:5-7}）。
在这些原电池里，铁是负极，碳是正极。
这时，作为负极的铁就失去电子而被氧化：
\[ \ce{ Fe - 2e = Fe^2+ } \]

在正极，溶液里的 \ce{H+} 得到电子而被还原成氢原子，氢原子又互相结合成氢分子，在碳的表面放出：
\begin{gather*}
  \ce{2H+ + 2e = 2H} \\
  \ce{2H = H_2 ^}
\end{gather*}

\begin{figure}
  \begin{minipage}[b]{0.48\linewidth}\centering
    \includegraphics{5-7.pdf}
    \caption{钢铁表面形成的微小原电池示意图}\label{fig:5-7}
  \end{minipage}
  \begin{minipage}[b]{0.48\linewidth}\centering
    \includegraphics{5-8.pdf}
    \caption{钢铁的电化腐蚀示意图}\label{fig:5-8}
  \end{minipage}
\end{figure}

钢铁在潮湿的空气里通过原电池反应而发生了电化腐蚀（\cref{fig:5-8}）。

上述腐蚀实际上是在酸性较强的溶液里进行的，在腐蚀过程里有氢气放出，这种腐蚀通常叫做析氢腐蚀。

在一般情况下，如果钢铁表面吸附的水膜酸性很弱或者呈中性，那么在负极上也是铁失去电子而被氧化成 \ce{Fe^2+}，而在正极上主要是溶解于水膜里的氧气得到电子而被还原：
\begin{gather*}
  \ce{2Fe - 4e = 2Fe^2+} \\
  \ce{2H2O + O2 + 4e = 4OH- }
\end{gather*}

所以，空气里的氧气溶解于水膜里也能促使钢铁腐蚀，这种腐蚀通常叫做吸氧腐蚀。
实际上，钢铁等金属的腐蚀主要是这种吸氧腐蚀。

金属的腐蚀除电化腐蚀外，还有金属跟接触到的物质（一般是非电解质）直接发生化学反应而引起的一种腐蚀，这种腐蚀叫做化学腐蚀。
例如，铁在高温下跟氧气直接反应，化工厂里的氯气等跟铁或其它金属直接反应而发生的腐蚀，都属于化学腐蚀。
这一类腐蚀的化学反应比较简单，仅仅是铁等金属跟氧化剂之间的氧化—还原反应。

从本质上看，电化腐蚀和化学腐蚀都是铁等金属原子失去电子而被氧化的过程，但是电化腐蚀过程里伴有电流产生，化学腐蚀过程里却没有。
在一般情况下，这两种腐蚀往往同时发生，只是电化腐蚀比化学腐蚀要普遍得多。

金属腐蚀的现象非常普遍，象金属制成的日用品、生产工具、机器部件、船壳等保养不好，就会腐蚀，从而造成大量金属的损耗。
至于因设备腐蚀损坏而引起停工减产、产品质量下降、污染环境、危害人体健康，甚至造成严重事故的损失，那就更无法估计了。
因此，了解金属腐蚀的原因，掌握防护的方法，是具有十分重要的意义的。

既然金属腐蚀主要是由于金属跟周围物质发生氧化-还原反应所引起，那么，金属的防护当然也必须从金属和周围物质两方面来考虑。
生产上常用的一些防护方法有：
\subsubsection{改变金属的内部组织结构}
例如，把铬、镍等加入普通钢里制成的不锈钢，就大大地增加了钢铁对各种侵蚀的抵抗力。

\subsubsection{在金属表面覆盖保护层}
在金属表面覆盖致密的保护层，从而使金属制品跟周围物质隔离开来，这是一种普遍采用的防护方法。

\begin{Reading}[补充阅读]{}
以钢铁的防护为例，根据保护层成分的不同，可分为如下几种：
\begin{enumerate}
  \item 在钢铁表面涂矿物性油脂、油漆或覆盖搪瓷、塑料等物质。例如船身、车厢、水桶等常涂油漆; 汽车外壳常喷漆; 枪炮、机器常涂矿物性油脂等。
  \item 用电镀、热镀、喷镀的方法，在钢铁表面镀上一层不易被腐蚀的金属，如锌、锡、铬、镍等。这些金属都能氧化而形成一层致密的氧化物薄膜，从而阻止水和空气等对钢铁的腐蚀。例如，常用的白铁皮就是在薄钢板的表面热镀上一层锌; 做罐头用的马口铁就是在薄钢板上热镀上一层无毒而又耐腐蚀的锡; 自行车钢圈就是用电镀的方法镀上一层既耐腐蚀又耐磨的铬或镍。
  \item 用化学方法使钢铁表面生成一层致密而稳定的氧化膜。例如，工业上常借一些溶液的氧化作用，在机器零件、精密仪器、枪炮等钢铁制件的表面上形成一层致密的黑色的四氧化三铁薄膜。
\end{enumerate}
\end{Reading}

\subsubsection{电化学保护法}
人们还可以利用原电池的化学原理来进行金属的防护，只要能够把引起金属发生电化腐蚀的原电池反应消除，金属的腐蚀自然就可以防止了。电化学保护法可分为阳极保护和阴极保护两大类。应用较多的是阴极保护。
\begin{Reading}[补充阅读]{}
\paragraph{外加电流的阴极保护法}\mbox{}\par
外加电流的阴极保护法的基本原理可简单地用\cref{fig:5-9} 表示。
这种方法是把要保护的钢铁设备作为阴极，另外用不溶性电极作为辅助阳极，两者都放在电解质溶液里，接上外加直流电源。
通电后，大量电子被强制流向被保护的钢铁设备，使钢铁表面产生负电荷（电子）的积累。
这样就抑制了钢铁发生失去电子的作用，从而防止了钢铁的腐蚀\footnote{金属氧化所生成的电子流跟外加电流的方向是相反的，只要外加足够强的电压，金属腐蚀而产生的原电池电流就不能被输送，因而腐蚀也就不能发生。根据这种防护原理，要保护的金属是在阴极。因此，这种方法叫做外加电流的阴极保护法。}。例如，对钢闸门采取涂料和外加电流的联合防腐的新方法，就可以收到显著的效果。

\begin{figurehere}
  \includegraphics{5-9.pdf}
  \caption{外加电流的阴极保护法示意图}\label{fig:5-9}
\end{figurehere}

\paragraph{牺牲阳极的阴极保护法}\mbox{}\par
牺牲阳极\footnote{这里的防护原理跟原电池的反应原理一样，负极上起的是氧化反应，通常叫做牺牲阳极。这种方法也由此得名。}的阴极保护法是在要保护的钢铁设备上联结一种更易失去电子的金属或合金。
例如，钢闸门的保护，有的就应用这种方法。
它是用一种比铁更为活泼的金属，如锌等，联结在钢闸门上。
这样，当发生电化腐蚀时，被腐蚀的是那种比铁活泼的金属，而铁被保护了（\cref{fig:5-10}）。
通常在轮船的尾部和在船壳的水线以下部分，装上一定数量的锌块，来防止船壳等的腐蚀，就是应用的这种方法。

\begin{figurehere}
  \includegraphics{5-10.pdf}
  \caption{牺牲阳极的阴极保护法示意图}\label{fig:5-10}
\end{figurehere}

目前，电化学保护法除应用于海水或河道中钢铁设备的保护外，还应用于防止电缆、石油管道、地下设备和化工设备等的腐蚀。

此外，还有缓蚀剂法，等等。
\end{Reading}

\begin{Practice}[习题]
\begin{question}
  \item 举例说明什么是原电池？它为什么能产生电流？
  \item 为什么实验室用锌跟稀硫酸反应制取氢气时，用粗锌比用纯锌反应快？为什么往硫酸中加入少量硫酸铜溶液能使制取氢气的反应加快？
  \item 为什么钢铁制品的表面保持清洁干燥就不容易生锈？为什么轮船在海水中行驶比在河水中行驶更易腐蚀？为什么钢铁制品的表面覆盖一层四氧化三铁可防止腐蚀，而表面覆盖一层铁锈（三氧化二铁的水合物）却不能防止腐蚀？
  \item 说明下列现象的原因：
  \begin{tasks}[label=(\arabic*)]
    \task 浸在水里的铁柱，与水面接触的部分比在水下的部分容易腐蚀；
    \task 用铝质铆钉来铆接铁板，铁板不易生锈，用铜质铆钉来铆接就容易生锈。
  \end{tasks}
  \item 镀层破损后，为什么镀锌钢板（白铁）比镀锡钢板（马口铁）耐腐蚀？
\end{question}
\end{Practice}

\section{电解和电镀}
\cref{sec:battery}我们学习了原电池，主要讨论化学能转变成电能的过程。
这一节要讨论电能怎样转变成化学能，以及电解原理的重要应用。

\subsection{电解的原理}
电解质溶液的导电跟金属的导电是不同的。
金属导电时，金属本身看不出变化。
电解质溶液导电时，电解质发生了明显变化。
\begin{Experiment}*[righthand ratio=0.6]
  如\cref{fig:5-11} 所示的 U 形管中,注入 \ce{CuCl2} 溶液，插入两根石墨棒作电极，用湿润的碘化钾淀粉试纸放在阳极碳棒附近，检验放出的气体。接通直流电源，观察管内发生的现象。
  \tcblower
  \begin{figurehere}
    \includegraphics{5-11.pdf}
    \caption{\ce{CuCl2} 溶液电解实验装置示意图}\label{fig:5-11}
  \end{figurehere}
\end{Experiment}
从实验可以看到，接通直流电源不久，在作为阴极的碳棒上有一层铜覆盖在它的表面，说明有铜析出。
在作为阳极的碳棒上，有气泡放出，从它的气味和它能使湿润的碘化钾淀粉试纸变蓝的特性，可以断定放出的气体是氯气。
由此可知，氯化铜溶液受到电流的作用，在导电的同时，发生了化学变化，生成了铜和氯气。

通电时，为什么氯化铜会分解成铜和氯气呢？

我们知道，氯化铜是强电解质，它在水里能电离成 \ce{Cu^2+} 和 \ce{Cl-}\footnote{实际上都是水合离子，为了简便起见写成简单离子。}：
\[ \ce{ CuCl2 = Cu^2+ + 2Cl- } \]

\begin{figure}
  \includegraphics{5-12.pdf}
  \caption{通电前后溶液里离子移动示意图}\label{fig:5-12}
\end{figure}
通电前，\ce{Cu^2+} 和 \ce{Cl-} 在水里自由地移动着；通电后,这些自由移动着的离子，在电场的作用下，改作定向移动。
根据异性相吸的原理，带负电的氯离子向阳极移动，带正电的铜离子向阴极移动，如\cref{fig:5-12} 所示，在阳极，氯离子失去电子而氧化成氯原子，然后两两结合成氯分子，从阳极放出。
在阴极，铜离子获得电子而还原成铜原子，就覆盖在阴极上。它们的反应可分别表示如下：
\begin{align*}
  \text{阳极}\quad \ce{2Cl- - 2e} &\ce{ = } \ce{2Cl}\text{ （氧化反应）}\\
  \ce{2Cl} &\ce{ = } \ce{Cl2 ^} \\
  \text{阴极}\quad \ce{Cu^2+ + 2e} &\ce{ = } \ce{Cu}\text{ （还原反应）}
\end{align*}


这种\emph{使电流通过电解质溶液而在阴阳两极引起氧化—还原反应的过程叫做}\Concept{电解}。
这种借助于电流引起氧化—还原反应的装置，也就是把电能转变为化学能的装置，叫做电解池或电解槽。
跟直流电源的负极相联的电极是电解池的阴极。
通电时，电子从电源的负极沿导线流入电解池的阴极。
跟直流电源的正极相联的电极是电解池的阳极。
通电时，电子从电解池的阳极流出，并沿导线流回电源的正极。
这样，电流就依靠溶液里阴阳离子的定向移动而通过溶液，所以电解质溶液的导电过程，就是电解质溶液的电解过程。
电解时，阳离子在阴极上得到电子发生还原反应; 阴离子在阳极上失去电子发生氧化反应。

氯化铜在水溶液里电解的化学方程式就是阳极上的反应和阴极上的反应的总和。
\[ \ce{ CU^2+ + 2Cl- \xlongequal{\text{电解}} Cu + Cl2 ^ }\]

在上面叙述氯化铜电解的过程中，没有提到溶液里的 \ce{H+} 和 \ce{OH-}，其实 \ce{H+} 和 \ce{OH-} 虽少，但确是存在的，只是它们没有参加电极反应。
也就是说在氯化铜溶液中，除 \ce{Cu^2+} 和 \ce{Cl-} 外，还有 \ce{H+} 和 \ce{OH-}，电解时，移向阴极的离子有 \ce{Cu^2+} 和 \ce{H+}，因为 \ce{Cu^2+} 比 \ce{H+} 容易得到电子，所以 \ce{Cu^2+} 在阴极上得到电子析出金属铜。
移向阳极的离子有 \ce{OH-} 和 \ce{Cl-}，在这样的实验条件下，\ce{Cl-} 比 \ce{OH-} 容易失去电子，所以 \ce{Cl-} 在阳极上失去电子，生成氯气。

\subsection{电解原理的应用}
\subsubsection{电解食盐水以制取氯气和烧碱}
工业上用电解饱和食盐水溶液的方法来制取烧碱、氯气和氢气，食盐水的电解原理同上面 \ce{CuCl2} 溶液的电解原理是相似的。

现在让我们仔细观察下面的实验。
\begin{Experiment}*[righthand ratio=0.6]
按\cref{fig:5-13} 装置，在 U 形管里倒入饱和食盐水，插入一根碳棒作阳极，一根铁棒作阴极。
同时在两边管中滴入几滴酚酞试液，并用湿润的碘化钾淀粉试纸检验阳极放出的气体。
接通直流电源后，注意管内发生的现象。
\tcblower
\begin{figurehere}
  \includegraphics{5-13.pdf}
  \caption{饱和食盐水电解实验装置}\label{fig:5-13}
\end{figurehere}
\end{Experiment}
从实验可以看到两极都有气体放出，阳极放出的气体，有刺激性气味，且能使湿润的碘化钾淀粉试纸变蓝，证明是氯气。
阴极放出的气体是氢气。
同时发现阴极附近溶液变红，说明溶液里有碱性物质生成。

这是因为食盐水里氯化钠完全电离，水分子也微弱电离，因而存在着 \ce{Na+}、\ce{H+}、\ce{Cl-}、\ce{OH-} 四种离子。
\begin{gather*}
  \ce{NaCl \xlongequal{\quad} Na+ + Cl-} \\
  \ce{H2O <=> H + + OH-}
\end{gather*}
当接通直流电源后，带负电的 \ce{OH-} 和 \ce{Cl-} 移向阳极，带正电的 \ce{Na+} 和 \ce{H+} 移向阴极。如\cref{fig:5-13} 所示。

在这样的电解条件下，在阳极，\ce{Cl-} 比 \ce{OH-} 容易失去电子被氧化成氯原子，氯原子两两结合成氯分子放出。
\begin{gather*}
  \ce{ 2Cl- - 2\mathrm{e} \xlongequal{\quad} 2Cl }\text{ （氧化反应）}\\
  \ce{ 2Cl \xlongequal{\quad} Cl2 ^}
\end{gather*}
在阴极，\ce{H+} 比 \ce{Na+} 容易得到电子，因而 \ce{H+} 不断从阴极获得电子被还原为氢原子，氢原子两两结合成氢分子从阴极放出。
\begin{gather*}
  \ce{ 2H+ + 2\mathrm{e} \xlongequal{\quad} 2H }\text{ （还原反应）}\\
  \ce{ 2H \xlongequal{\quad} H2 ^}
\end{gather*}

由于 \({\mathrm{H}}^{ + }\) 在阴极上不断得到电子而生成氢气放出,破坏了附近的水的电离平衡,水分子继续电离成 \({\mathrm{H}}^{ + }\) 和 \({\mathrm{{OH}}}^{ - },{\mathrm{H}}^{ + }\) 又不断得到电子,结果溶液里 \({\mathrm{{OH}}}^{ - }\) 的数目相对地增多了。因而阴极附近形成了氢氧化钠的溶液。电解饱和食盐水的总的化学方程式可以表示如下：
\[ \ce{ 2NaCl + 2H2O \xlongequal{\text{电解}} 2NaOH + H2 ^ + Cl2 ^ } \]

\par\medskip\noindent
\begin{minipage}{0.5\linewidth}\parindent2em
目前我国工业上大多采用立式隔膜电解槽，如\cref{fig:5-14} 所示。这种隔膜电解槽的阳极是由金属钛或石墨制成，阴极是铁丝网作成的，网上吸附着一层石棉绒作为隔膜，这层隔膜把电解槽隔成阳极室和阴极室。
隔膜能阻止气体分子穿过，但不能阻止水分子和离子穿过。
这样既能防止阴极产生的氢气和阳极产生的氯气相混和而引起爆炸，又能避免氯气和氢氧化钠作用生成次氯酸钠，而影响烧碱的产量和质量。
氯气跟氢氧化钠的反应可以表示如下：
\end{minipage}\hfill
\begin{minipage}{0.45\linewidth}
\begin{figurehere}
  \includegraphics{5-14.pdf}
  \caption{立式隔膜电解槽示意图}\label{fig:5-14}
\end{figurehere}
\end{minipage}

\[ \ce{ 2NaOH + Cl2 \xlongequal{\quad} NaClO + NaCl + H2O }\]

原料食盐水要经过净制，以除去里面的 \ce{Mg^2+}、\ce{Ca^2+}、\ce{SO4^2-} 等离子，以免成品混入杂质。
同时也可以防止电解过程中产生的 \ce{Mg(OH)2} 之类的不溶性杂质，堵塞隔膜孔隙。
经过净制的浓盐水不断由导管缓缓进入阳极室，通过隔膜进入阴极室。
氯气由阳极室上部的管子放出，氢气由阴极室上部的管子放出。含有 \ce{NaOH} 和 \ce{NaCl}（未电解的食盐）的电解液从阴极室底部的碱液出口管导出。
最后加热蒸发流出的混和溶液，溶解度较小的 \ce{NaCl} 大部分结成晶体析出。
除去 \ce{NaCl} 后的浓 \ce{NaOH} 溶液（即液碱），可以直接供工业上使用，也可以进一步精制，以制得固态的 \ce{NaOH}。

\subsubsection{电镀}
\Concept{电镀}\emph{是应用电解原理在某些金属表面镀上一薄层其它金属或合金的过程}。
电镀的目的主要是使金属增强抗腐蚀能力、增加美观和表面硬度。
镀层金属通常是一些在空气或溶液里不易起变化的金属（如铬、锌、镍、银）和合金（铜锡合金、铜锌合金等）。

电镀时，把待镀的金属制品作阴极，把镀层金属作阳极，用含有镀层金属的离子的溶液作电镀液。
在直流电的作用下，镀件表面就覆盖上一层均匀光洁而致密的镀层。
现以镀锌为例来说明电镀的过程：
\begin{Experiment}*[righthand ratio=0.5]\label{expr:5-6}
  按\cref{fig:5-15} 装置，在大烧杯里加入锌电镀液（含 \ce{ZnCl2} 的溶液）。用锌片作阳极，把待镀的铁制品（即镀件）作阴极。接通直流电源几分钟后，就可看到镀件的表面被镀上了一层锌。
  \tcblower 
  \begin{figurehere}
    \includegraphics{5-15.pdf}
    \caption{电镀锌的实验装置}\label{fig:5-15}
  \end{figurehere}
\end{Experiment}

上述镀锌主要过程可以表示如下：

通电前：
\[ \ce{ ZnCl2 = Zn^2+ + 2Cl- }\]

通电后：
\begin{align*}
  \text{在阴极}\qquad & \ce{ Zn^2+ + 2e =  Zn }\ \text{（还原反应）}\\
  \text{在阳极}\qquad & \ce{ Zn - 2e = Zn^2+ }\ \text{（氧化反应）}
\end{align*}

除 \ce{Zn^2+}、\ce{Cl-} 两种离子外,电镀液中还有由 \ce{H2O} 电离出的 \ce{H+} 和 \ce{OH-}，在电镀所控制的条件下，这些离子一般不起反应。

电镀的结果，阳极的锌（镀层金属）不断减少，阴极的锌（镀件上）不断增加，减少和增加的锌量相等。因此溶液里的 \ce{ZnCl2} 含量是保持不变的。

由此可见，镀锌过程包括了在阴极 \ce{Zn^2+} 得到电子和在阳极 \ce{Zn} 失去电子的氧化—还原反应的过程。因此，电镀过程实质上就是一个电解过程。它的特点是阳极本身也参加了电极反应（失去电子而溶解）。

为了使镀层光滑牢固，工业生产上电镀液的配方往往很复杂，而且过去镀银或镀别的金属，都常用氰化物来配制电镀液。
由于 \ce{CN-} 有剧毒，对电镀工人的健康会造成损害，排出含 \ce{CN-} 的污水、废气会严重污染环境，所以，现在正在积极研究和逐步使用无氰电镀液。\cref{expr:5-6} 使用的就是无氰电镀液。

除了可以对金属制品进行电镀外，还可以对塑料进行电镀。
由于塑料是非导体，不能直接进行电镀。
在镀前要对塑料表面进行预处理，除去塑料表面的油和杂质，使塑料表面洁净，再沉积一层导电的金属膜，然后才能跟金属电镀一样把它作为阴极，进行电镀。
塑料制品经过电镀金属层后，具有分量轻、能导电、外表美观等优点。
不仅能代替金属铜和铝，还能减少加工工序，降低成本。
因此，目前塑料电镀工艺已应用于电子、光学仪器、机床按钮、轻工产品等各方面。

电解原理的应用还有电冶、电解精炼等等，这些将在以后有关内容里学到。

\begin{Practice}[习题]
\begin{question}
  \item 什么叫做电解？电解和电镀有什么相同和不同的地方？
  \item 试说明电解水（加入少量 \ce{H2SO4} 或 \ce{NaOH} 溶液）时，电解槽的阴、阳极分别放出氢气和氧气的原因。写出化学反应的方程式。
  \item 电镀有哪些优点，在工业上有哪些应用？
  \item 为什么电解食盐溶液的电解槽要设置隔膜？为什么电解食盐溶液能制得烧碱、氯气和氢气？用化学方程式表示食盐溶液的电解过程。
  \item 电解硫酸铜溶液时，在电解槽阴、阳两极各发生什么化学变化？写出反应的化学方程式。
  \item 如果我们要在铜制品表面镀上一层银，电解槽应用什么做阴、阳极？电镀液中必须包含什么盐类？
\end{question}
\end{Practice}


\section*{内容提要}
\addcontentsline{toc}{section}{内容提要}\setcounter{subsection}{0}
\subsection*{一、强电解质和弱电解质}
强电解质在水溶液里全部电离为离子，如强酸、强碱和大部分盐类都是强电解质。

弱电解质在水溶液里只有部分电离为离子，溶液里还存在未电离的电解质分子，如弱酸、弱碱等都是弱电解质。
\subsection*{二、弱电解质的电离}
\begin{enumerate}[1.]
  \item 弱电解质的电离平衡\quad 弱电解质的电离是一个可逆过程。溶液中未电离的电解质分子和电解质部分电离生成的离子处于平衡状态，如醋酸的电离：
  \[ \ce{CH3COOH <=> H+ + CH3COO-} \]

  这种电解质溶液中电解质分子电离成离子的速度跟离子重新结合成分子的速度相等时，叫做电离平衡。

  电离平衡是动态平衡，当条件（浓度、温度）发生变化时，平衡就向能够使这种变化减弱的方向移动。
  \item 电离度
  \[ \text{电离度}(\alpha)=\frac{\text{已电离的电解质分子数}}{\text{溶液中原有电解质分子的总数}}\times 100\%\]
  \item 电离常数
  
  当弱电解质 \ce{AB} 在水溶液里电离：
  \begin{gather*} 
    \ce{AB <=> A+ + B-}\\
    K_\text{电离}=\frac{[\ce{A+}][\ce{B-}]}{[\ce{AB}]}
  \end{gather*}
  式中 $[\ce{A+}]$、$[\ce{B-}]$、$[\ce{AB}]$ 分别表示 \ce{A+} 离子、\ce{B-} 离子和 \ce{AB} 分子的摩尔浓度，$K_\text{电离}$ 为 \ce{AB} 溶液的电离常数。

  电离常数的大小，跟电离度一样能衡量电解质的相对强弱。不过电离度跟电解质溶液的浓度有关，而电离常数跟浓度无关。
\end{enumerate}

\subsection*{三、水的离子积和溶液的 pH 值}
\begin{enumerate}[1.]
  \item 水的离子积\quad 在常温时，水中 $[\ce{H+}]\times[\ce{OH-}]=\num{1e-14}$，这个常数叫做水的离子积。
  \item 溶液的 pH 值
  
  水溶液里 \ce{H+} 浓度的负对数叫做溶液的 pH 值。即：
  \[ \text{pH} = -\lg [\ce{H+}]\]
\end{enumerate}

\subsection*{四、盐类的水解}
盐类的离子跟水电离出来的 \ce{H+} 或 \ce{OH-} 相结合生成弱电解质的反应，叫做盐类的水解反应。它是中和反应的逆反应。 例如：
\[ \ce{ CH3COONa + H2O } \xLongleftrightarrow[\text{中和}]{\text{水解}} \ce{NaOH + CH3COOH }\]

可写成离子方程式：
\[ \ce{ CH3COO- + H2O } \xLongleftrightarrow[\text{中和}]{\text{水解}} \ce{CH3COOH + OH-}\]

各种类型盐的水解情况比较：

\begin{tablehere}
  \begin{tblr}{colspec={cX[c]cX[c]},hline{2}=0.8pt}
    盐的类型   & 举例 & 水解情况 & 溶液的酸碱性 \\
    强酸弱碱盐 & \ce{NH4Cl}、\ce{Al2(SO4)3}、\ce{FeCl3} & 水解 & 酸性，$\mathrm{pH}<7$ \\
    强碱弱酸盐 & \ce{CH3COONa}、\ce{Na2CO3}、\ce{NaHCO3} & 水解 & 碱性，$\mathrm{pH}>7$ \\
    弱酸弱碱盐 & \ce{CH3COONH4}、\ce{NH4CN}、\ce{(NH4)2S} & 水解 & 中性、酸性或碱性（随酸碱的相对强弱而异） \\
    强酸强碱盐 & \ce{NaCl}、\ce{KNO3}、\ce{Na2SO4} & 不发生水解 & 中性，$\mathrm{pH}=7$ \\
  \end{tblr}
\end{tablehere}


\subsection*{五、酸碱的当量浓度和中和滴定}
\begin{enumerate}[1.,itemsep=7pt]
  \item $\text{酸的克当量}=\dfrac{\qty{1}{mol} \text{酸的质量}(\unit{g}) }{\qty{1}{mol} \text{酸所提供}\ce{H+}\text{的摩尔数}}$
  \item $\text{碱的克当量}=\dfrac{\qty{1}{mol} \text{碱的质量}(\unit{g}) }{\qty{1}{mol} \text{碱所提供}\ce{OH-}\text{的摩尔数}}$
  \item 当量浓度\quad 溶液的浓度用每升溶液里所含溶质的克当量数表示的浓度。常用 $N$ 表示。
  \item 酸跟碱以相等的克当量数发生反应。
  \item 中和滴定\quad 根据 $N_\text{酸}V_\text{酸}=N_\text{碱}V_\text{碱}$ 的公式，当加已知浓度的酸（或碱）溶液于未知浓度的碱（或酸）溶液时，可测知未知碱 （或酸）溶液的浓度，这种方法叫做中和滴定。
  \item 中和热\quad 在稀溶液中，酸跟碱发生中和反应而生成 \qty{1}{mol} 水，这时的反应热就是中和热。
\end{enumerate}

\subsection*{六、原电池、电解和电镀}
原电池、电解和电镀的反应都是在电极上发生的氧化—还原反应。
原电池是把化学能转变为电能的装置，金属的电化腐蚀是通过原电池反应而引起的。
电解和电镀都是把电能转变为化学能的过程。
原电池和电解池的情况对比如下：

\begin{table}
  \begin{tblr}{colspec={cX[2,l]X[2,l]X[2,l]X[l]},hline{2}=0.8pt}
    & 电极名称 & 电流方向 & 电极反应 & 能量转变 \\
    {原\\电\\池} & 正、负极由电极本身决定：相对活泼金属为负极，较不活泼金属（根据金属活动性顺序）为正极。& 电子由负极移向正极。电流由正极流向负极。& 负极：失去电子，起氧化反应。正极：阳离子得到电子，起还原反应。& 由化学能转变为电能。 \\
    {电\\解\\池} & 阴、阳两极由外电源决定：跟直流电源负极相连的电极为阴极；跟正极相连的电极为阳极。& 电子由直流电源的负极流出，经过导线到达电解池的阴极，然后由离子通过电解液到达阳极，离子在阳极失去电子，电子再回到直流电源的正极 & 阴极：较易得电子的阳离子优先 在阴极得到电子而起还原反应。阳极（惰性）：较易失电子的阴离子优先在阳极失去电子而起氧化反应。& 由电能转变为化学能。\\
  \end{tblr}
\end{table}

\begin{Review}
\begin{question}
  \item 在稀释弱电解质浓溶液时，一般为什么开始时溶液的导电性逐渐增强，以后又逐渐减弱？
  \item 在纯水里加入少量的酸或碱后，水的离子积有无变化？ 水中 \ce{H+} 浓度有无变化？
  \item 当某溶液中 \ce{H+} 浓度为 $0.2M$ 时，它的 pH 值是多少？
  \item 下列溶液中哪一种溶液的 $[\ce{H+}]$ 最大？
  \begin{tasks}(2)
    \task \qty{30}{ml} $0.1M$ 醋酸；
    \task \qty{20}{ml} $0.1M$ 盐酸；
    \task \qty{10}{ml} $0.1M$ 硫酸。
  \end{tasks}
  \item 回答下列问题：
  \begin{tasks}[label=(\arabic*)]
    \task 两种浓度不同的氢氧化钠溶液的 pH 值的差为 1，求它们 $[\ce{OH-}]$ 的比值。
    \task 现有盐酸的 pH 值为 4，如果要配制 pH 值为 5 的盐酸应加几倍体积的水冲稀？
  \end{tasks}
  \item 一种溶液里可能含有 \ce{OH-}、\ce{Cl-}、\ce{NO3-}、\ce{CO3^2-} 和 \ce{SO4^2-} 五种阴离子中的部分离子，取少量这种溶液分盛于三个试管里，进行如下实验：
  \begin{tasks}[label=(\arabic*)]
    \task 向第一个试管里滴入酚酞溶液时变红，加热浓缩后，加浓 \ce{H2SO4} 和铜屑，再加热产生红棕色气体；
    \task 向第二个试管里加入 \ce{Ba(NO3)2} 溶液生成白色沉淀；
    \task 向第三个试管里逐渐加入稀 \ce{HNO3} 生成无色气体，该气体能使澄清石灰水变浑浊，继续加稀 \ce{HNO3} 于溶液里使溶液显酸性时，再加入 \ce{Ba(NO3)2} 溶液不产生沉淀，加入 \ce{AgNO3} 溶液也不产生沉淀。

    从上述实验试判断这种溶液里有哪几种离子？说明理由，写出反应的离子方程式。
  \end{tasks}
  \item 把 \qty{50}{ml} $0.15N$ 盐酸跟 \qty{25}{ml} $0.15N$ \ce{NaOH} 溶液混和，求该溶液的 pH 值（假定这两种溶液混和后，其体积为原来体积的和）。
  \item 硫酸铵、硝酸铵和氯化铵都是白色晶体，如何把它们区别开？写出离子方程式。
  \item 氯化钾中含有少量硫酸钾的杂质，怎样除去？写出有关的离子方程式。
  \item \qty{25}{\celsius} 时，把 \qty{0.5}{M} 氨水加水稀释到原浓度的 1/10，溶液中 \ce{OH-} 浓度有什么变化？电离度有什么变化？
  \item 计算下列各溶液 pH 值的大小。
  \begin{tasks}[label=(\arabic*)]
    \task \qty{20}{ml} $0.1M$ \ce{HCl} 和 \qty{20}{ml} $0.1M$ \ce{NaOH} 溶液混和，
    \task \qty{20}{ml} $0.1M$ \ce{HCl} 和 \qty{30}{ml} $0.1M$ \ce{NaOH} 溶液混和，
    \task \qty{30}{ml} $0.1M$ \ce{HCl} 和 \qty{20}{ml} $0.1M$ \ce{NaOH} 溶液混和。
  \end{tasks}
  \item 在摩尔浓度相同的 \ce{CH3COOH} 和 \ce{CH3COONa} 的溶液里，哪个溶液里 \ce{CH3COO-} 浓度较大？ 为什么？
  \item 分别绘图说明原电池的正极、负极以及电子移动方向，和电解池的阴极、阳极以及电子移动的方向。
  \item 在电解和电镀时，阳极发生什么反应，阴极发生什么反应？用氧化—还原的观点加以说明。
  \item 在电镀过程中，从理论上说电解质溶液的浓度有没有变化？为什么？对于电解过程是否也能得到同样的结论？举例说明。
\end{question}
\end{Review}